\documentclass[10pt,a4paper]{amsart}
\usepackage[margin=19mm]{geometry}
\usepackage{amsmath,amssymb,mathtools}
\usepackage[scr=rsfs,cal=euler]{mathalfa}
\usepackage{xfrac}
\usepackage[unicode,hidelinks,bookmarks=false]{hyperref}
\newcommand{\ie}{i.e.\ }
\newcommand{\cf}{cf.\ }
\newcommand{\resp}{respectively\ }
\newcommand{\Subsection}{\subsection}
\providecommand{\nobreakdash}{\nobreak\hbox{-}}
\setlength{\emergencystretch}{2em}
\multlinegap0pt
	\newcommand{\zed}{\ensuremath{\mathbb{Z}}}
	\newcommand{\kew}{\ensuremath{\mathbb{Q}}}
	\newcommand{\oh}{\ensuremath{\mathcal{O}}}
\usepackage{amssymb}
\usepackage{enumerate}
\usepackage{csquotes}
\usepackage{mathtools}
\newtheorem{lemma}{Lemma}
\title{$p$-adic Equidistribution of Special Loci in a Product of Modular Curves}
\author{Dan Townsend}
\date{Source: arXiv:2608.23052v1 (CC BY 4.0)}
\begin{document}
\maketitle

\noindent\emph{Source and licence.} $p$-adic Equidistribution of Special Loci in a Product of Modular Curves. Version arXiv:2608.23052v1 (CC BY 4.0), distributed by arXiv under the Creative Commons Attribution 4.0 International licence, http://creativecommons.org/licenses/by/4.0/, as recorded in the arXiv licence metadata for this exact version and shown on the abstract page https://arxiv.org/abs/2608.23052v1. Full licence notice: \texttt{licenses/arxiv-2608.23052-CC-BY-4.0.txt}.\par\smallskip

\noindent\textit{Editorial scope.} Counted unit: the article's lemma approxlemma (source0.tex lines 538--541). The article's complete original proof of it is delivered with it (source0.tex lines 543--567, one \texttt{proof} environment of 25 lines). No mathematics is written, repaired or re-derived: the statement and the proof are the article's own lines, in the article's own notation, and no retained line is substituted or re-flowed.  No further source-local context is needed: every cross-reference of every retained line resolves inside the two delivered spans.  Nothing was omitted from any retained span, so the package carries no omission record.  No retained line contains a citation, so the wrapper carries no bibliography and no bibliography entry.  No retained line contains a figure environment, an image-inclusion command or drawing code, so no image is delivered and the package declares no asset dependency.  Editorial formatting only, in addition: an AMS \texttt{amsart} preamble that restates the article's own declarations and macros used by the retained lines, namely the environments \texttt{lemma} on separate counters, together with the standard packages those lines need.  The retained environments print in this document as Lemma 1 in order of appearance, whereas the article prints the same items with the numbering of its own full document.  The complete native source of the article as distributed in the arXiv e-print for this exact version is bundled unchanged as \texttt{sources/208386/source0.tex}.
	\begin{lemma} \label{approxlemma}
		Let $\oh$ be an order of discriminant prime to $p$ in the ring of integers of the quadratic imaginary number field $\kew(\sqrt{-d})$, with $d>0$ a squarefree integer. Let $I\triangleleft \oh$ be a non-zero ideal. Then the collection of elements
		\[\{ \frac{\bar{z}}{z}\ |\ z\in I,\ v_p(z)\leq v_p(\bar{z})\}\] is dense in $\zed_p$. 
	\end{lemma}

	\begin{proof}
		
		Write $\oh = \zed\oplus e\omega \zed$ where $\{1,\omega\}$ form an integral basis for the ring of integers in $\kew(\sqrt{-d})$. The value of Disc$(\oh)$ is either $-e^2d$ or $-4e^2d$, so the hypotheses of the Lemma imply that $d$ and $e$ are coprime to $p$. Choose a non-zero element $\beta+\gamma\sqrt{-d}\in I$. After perhaps multiplying by 2, may assume $\beta,\gamma\in \zed$. 
		
		
	
	All elements of the form $(u+v\sqrt{-d})(\beta+\gamma\sqrt{-d})$, where $u,v\in\zed$ and $e|v$, all lie in $I$ (although of course not every element need have this form). Thus the Lemma would follow from finding, for every $t\in \zed_p$ and for every sufficiently large $r>0$, a solution $u\in \zed$, $v\in e\zed$, to the congruence 
	\begin{align}\label{Case1Cong}
		t\equiv \frac{u\beta-vd\gamma-(v\beta+u\gamma)\sqrt{-d}}{u\beta-vd\gamma+(v\beta+u\gamma)\sqrt{-d}}\pmod{p^r},
	\end{align}
	The congruence is equivalent to 	
		\[
		(1-t)(u\beta-vd\gamma) \equiv (1+t)(v\beta+u\gamma)\sqrt{-d} \pmod{p^r},\] as long as both sides are not identically zero.	Rearranging further, the above becomes:		
		\begin{align*} v(\beta\sqrt{-d}(1+t)+d\gamma(1-t))\equiv u(\beta(1-t)-\gamma\sqrt{-d}(1+t))\pmod{p^r}.\end{align*}		
		Write $A = \beta\sqrt{-d}(1+t)+d\gamma(1-t)$ and $B= \beta(1-t)-\gamma\sqrt{-d}(1+t)$. If $A=B=0$, then computing the ratio $\frac{1+t}{1-t}$ (or its inverse if $t=1$) from each of $A=0$ and $B=0$ gives $\sqrt{-d}(\beta^2+d\gamma^2)=0$, a contradiction to having chosen a non-zero element of $I$ in the first place. Now, fix $r>0$.
		
	
		
		\textbf{Case 1:} Suppose first that $R = v_p(A) \leq v_p(B)$ (so $R<\infty$). Further, set \[S = v_p(N(\beta+\gamma\sqrt{-d}))<\infty.\] Then choose $(u,v)\equiv(1,\frac{B}{A})\pmod{p^{r}}$, where this makes sense because $\frac{B}{A}\in \zed_p$. Then consider \begin{align*} z &= (u+v\sqrt{-d})(\beta+\gamma\sqrt{-d})  \\ &\equiv \frac{1}{A}(A+B\sqrt{-d})(\beta+\gamma\sqrt{-d}) \pmod{p^r} \\ &\equiv \frac{2\sqrt{-d}(\beta-\gamma\sqrt{-d})(\beta+\gamma\sqrt{-d})}{A} \pmod{p^r}  \\ &\equiv \frac{2\sqrt{-d}N(\beta+\gamma\sqrt{-d})}{A}\pmod{p^r}.\end{align*} Thus if  $r>S$, $v_p(z) = S-R$. Further, \begin{align*} \bar{z} &= (u-v\sqrt{-d})(\beta-\gamma\sqrt{-d})  \\ &\equiv \frac{1}{A}(A-B\sqrt{-d})(\beta-\gamma\sqrt{-d}) \pmod{p^r} \\ &\equiv \frac{2t\sqrt{-d}(\beta+\gamma\sqrt{-d})(\beta-\gamma\sqrt{-d})}{A} \pmod{p^r} \\ &\equiv \frac{2t\sqrt{-d}N(\beta+\gamma\sqrt{-d})}{A}\pmod{p^r}.\end{align*} So $v_p(\bar{z}) \geq S-R$ (with equality if $2t$ is a unit). Then $\frac{\bar{z}}{z}\equiv t\pmod{p^r}$ as desired. Since $(e,p^r)=1$, it is possible to do this while simultaneously forcing $v\equiv 0\pmod{e}$ by the Chinese Remainder Theorem. 
		
		\textbf{Case 2:} If instead $v_p(B)\leq v_p(A)$, take $(u,v) \equiv (\frac{A}{B},1)\pmod{p^r}$ and $z = (u+v\sqrt{-d})(\beta+\gamma\sqrt{-d})$. Now compute
		 \[
		 z \equiv \frac{1}{B}(A+B\sqrt{-d})(\beta+\gamma\sqrt{-d}) \equiv \frac{2\sqrt{-d}N(\beta+\gamma\sqrt{-d})}{B}\pmod{p^r},
		\] so $v_p(z)=S-R$ and \[\bar{z}\equiv \frac{1}{B}(A-B\sqrt{-d})(\beta-\gamma\sqrt{-d}) \equiv \frac{2t\sqrt{-d}N(\beta+\gamma\sqrt{-d})}{B}\pmod{p^r},\] whence $v_p(\bar{z})\geq S-R$, and $\frac{\bar{z}}{z}\equiv t \pmod{p^r}$ with $v_p(z)\leq v_p(\bar{z})$. Again, by the Chinese Remainder Theorem, it is possible to simultaneously force $v\equiv 0\pmod{e}$. 
		\end{proof}

\end{document}
